% Generated by tools/make_interview_tracks.py. Do not edit.
\section*{Quels chapitres travailler}
Chaque case compte les questions d'un chapitre étiquetées pour ce métier ou ce type de firme. QR désigne le chercheur quantitatif (quant researcher), QT le trader quantitatif, DS le data scientist, QD le développeur quantitatif, Dev le développeur logiciel, LLE l'ingénieur basse latence, Risque le quant de risque~; TM le teneur de marché, Cr la firme crypto, Bq la banque, HF le hedge fund. Les épreuves et les exercices de vitesse sans étiquette sont pour tout le monde.
\par\medskip\noindent{\footnotesize\setlength{\tabcolsep}{2.5pt}%
\begin{tabular}{@{}>{\raggedright\arraybackslash}p{6.2cm}rrrrrrrrrrrrr@{}}
\hline
Chapitre & QR & QT & Trader & Analyste & DS & QD & Dev & LLE & Risque & TM & Cr & Bq & HF \\
\hline
\hyperref[ch:iv3:algorithms]{Algorithmes} & 34 & 12 &  &  & 13 & 80 & 76 & 32 &  & 78 & 74 & 64 & 78 \\
\hyperref[ch:iv3:python-and-pandas]{Python et pandas} & 28 & 6 &  & 4 & 26 & 39 & 19 & 6 & 2 & 39 & 35 & 36 & 39 \\
\hyperref[ch:iv3:cpp]{C++} & 5 &  &  &  &  & 60 & 60 & 59 &  & 60 & 60 & 53 & 60 \\
\hyperref[ch:iv3:low-latency]{Faible latence} & 3 & 2 &  &  &  & 24 & 35 & 35 &  & 35 & 19 & 5 & 21 \\
\hyperref[ch:iv3:operating-systems-concurrency-and-networks]{Systèmes d'exploitation, concurrence et réseaux} &  &  &  &  &  & 40 & 40 & 35 &  & 40 & 40 & 27 & 40 \\
\hyperref[ch:iv3:system-design]{Conception de systèmes} & 6 & 2 &  & 1 & 7 & 40 & 25 & 20 & 4 & 36 & 38 & 30 & 36 \\
\hyperref[ch:iv3:sql-and-data]{SQL et données} & 12 & 2 &  &  & 25 & 25 & 25 & 2 &  & 24 & 24 & 20 & 24 \\
\hyperref[ch:iv3:smart-contracts-and-crypto-engineering]{Smart contracts et ingénierie crypto} &  & 7 &  &  &  & 30 & 26 & 2 & 4 &  & 30 &  &  \\
\hyperref[ch:iv3:reading-debugging-and-software-design]{Lecture, débogage et conception logicielle} & 9 & 4 &  &  & 7 & 40 & 40 & 24 &  & 34 & 26 & 29 & 31 \\
\hline
\end{tabular}}
